\chapter{Data Loading and Inspection}
\label{cp:dli}

\section{Initial Data Structure}
The initial inspection revealed a dataset of $48,842$ entries, with 6 numerical columns and 9 categorical (object-type) columns. Notably, the initial data check showed zero non-null values across all columns.


\vspace{1 \baselineskip}
\begin{center}
\begin{tabular}{||c | c||} 
 \hline
 Qualitative variables & Quantitative variables \\ [1 ex] 
 \hline\hline
 Workclass & Age \\ [1ex] 
 \hline
 Marital-status & Fnlwgt \\ [1ex] 
 \hline
 Education & Educational-num \\ [1ex] 
 \hline
 Occupation & Capital-gain  \\ [1ex] 
 \hline
 Relationship & Capital-loss  \\ [1ex] 
 \hline
 Race & Hours-per-week  \\ [1ex] 
 \hline
 Gender & -  \\ [1ex] 
 \hline
 Native-Country & -  \\ [1ex] 
 \hline
 Income (target) & -  \\ [1ex] 
 \hline
\end{tabular}
\end{center}


\vspace{1 \baselineskip}
\section{Handling Missing Values}
\label{sec:handling_missing}
The absence of null values was highly unusual and indicated that missing data was likely encoded with a special character, rather than the standard $\text{NaN}$ format. An inspection of the \texttt{workclass} values confirmed this, revealing the character '?'.

To rectify this, we replaced all occurrences of the '?' character with $\text{np.nan}$ across the entire dataset. A subsequent count of true missing values revealed the following:

\begin{itemize}
    \item \texttt{workclass}: $2,799$ missing values
    \item \texttt{occupation}: $2,809$ missing values
    \item \texttt{native-country}: $857$ missing values
\end{itemize}
This represents a significant loss of information ($\approx 5.7\%$ of rows contain missing data in \texttt{workclass} or \texttt{occupation}). Therefore, we assumed that simply deleting these rows was a bad strategy. Our next $\text{EDA}$ phase was tasked with determining the most appropriate imputation strategy for these $\text{NaN}$ values.

\vspace{2 \baselineskip}